\begin{abstract}
Large language models are absorbing more workplace tasks, increasingly handed to them by users who cannot fully judge or direct what comes back. When the output needs improving, those users do the only thing they can: ask the model to make it better, without saying what is wrong. We show that this undirected revision is where the collaboration quietly breaks down. Across 720 five-turn conversations spanning six models and forty tasks, models asked to revise without direction mostly decline to revise at all, dressing refusal as engagement: genuine revision falls from 39\% of responses at the second turn to 13\% by the fifth. When models do revise, they degrade their own work, quality falls 0.74 levels on a five-point scale ($p = 1.01 \times 10^{-4}$), and 62\% of output tokens are spent past the point of peak quality. The problem is direction, not capacity: a single targeted critique raises quality 1.16 levels and reverses the decline. Measuring any of this demands care, models wrap revisions in meta-commentary that inflates an automated judge's preference for the first draft from 56\% to 92\%, an effect that nearly vanishes once the commentary is stripped. As models take on work their users cannot direct, the remedy is not more iteration but evaluation: tell the model what is wrong, or do not ask it to revise.
\end{abstract}
